\section{Regime-Local Hybrid ROM Specialists}
\label{sec:regime_local_specialists}
\label{sec:hprs_specialists}

Each regime-local chart $\mathcal{C}^{\mathrm{ROM}}_r$ is augmented with a
complete autonomous reduced dynamical model, referred to as a
\emph{regime-local specialist}.  The specialists share a common
physics--data architecture but are parameterized independently on their
respective regime-local data.  Within a specialist, the term \emph{expert}
refers only to a neural component of the internal sparse mixture and should
not be confused with the specialist itself.  The construction has two
levels: a projected reduced model provides the local physical backbone, and
a structured sparse mixture of experts represents unresolved contributions
conditioned on the recent reduced trajectory.

\subsection{Semi-explicit hybrid reduced dynamics}
\label{subsec:semi_explicit_hybrid_dynamics}

For regime $r\in\mathcal{R}$, the velocity coefficients satisfy
\begin{equation}
  \dot{\boldsymbol{a}}^{(r)}
  =
  \mathscr{F}^{\mathrm{G}}_r
  \bigl(\boldsymbol{a}^{(r)},\boldsymbol{b}^{(r)};\mu\bigr)
  +
  \boldsymbol{\rho}^{u}_{\theta_r}
  \bigl(\boldsymbol{\xi}^{(r)}\bigr),
  \label{eq:hybrid_velocity_rhs}
\end{equation}
where $\mathscr{F}^{\mathrm{G}}_r$ is the chart-local Galerkin vector field
from~\eqref{eq:chart_local_galerkin_rhs_definition}.  The learned term
$\boldsymbol{\rho}^{u}_{\theta_r}$ therefore acts as a closure of the
projected dynamics rather than replacing them.

Pressure is treated as an algebraic variable.  After advancing the velocity,
the reduced Pressure--Poisson map provides the physics-based candidate
\begin{equation}
  \boldsymbol{b}^{PP,(r)}_{n+1}
  =
  \mathscr{Q}^{PP}_{r}
  \bigl(\boldsymbol{a}^{(r)}_{n+1};\mu\bigr).
  \label{eq:pressure_poisson_base}
\end{equation}
A data-driven candidate and a componentwise gate are evaluated from the same
specialist feature vector,
\begin{align}
  \widehat{\boldsymbol{b}}^{D,(r)}_{n+1}
  &=
  \boldsymbol{\rho}^{p}_{\theta_r}
  \bigl(\boldsymbol{\xi}^{(r)}_n\bigr),
  \label{eq:data_driven_pressure_candidate}\\
  \boldsymbol{\alpha}^{(r)}_n
  &=
  \operatorname{sigmoid}
  \left(
    \boldsymbol{\eta}^{p}_{\theta_r}
    \bigl(\boldsymbol{\xi}^{(r)}_n\bigr)
  \right),
  \qquad
  \boldsymbol{\alpha}^{(r)}_n\in(0,1)^{d_p}.
  \label{eq:adaptive_pressure_gate}
\end{align}
The pressure coefficients are then reconstructed as
\begin{equation}
  \boldsymbol{b}^{(r)}_{n+1}
  =
  \boldsymbol{\alpha}^{(r)}_n\odot
  \boldsymbol{b}^{PP,(r)}_{n+1}
  +
  \bigl(\boldsymbol{1}-\boldsymbol{\alpha}^{(r)}_n\bigr)
  \odot
  \widehat{\boldsymbol{b}}^{D,(r)}_{n+1}.
  \label{eq:adaptive_pressure_closure}
\end{equation}
Equations~\eqref{eq:hybrid_velocity_rhs}--\eqref{eq:adaptive_pressure_closure}
form a partitioned semi-explicit ROM: velocity is advanced by a differential
model, whereas pressure is recovered from an algebraic relation conditioned
on the updated velocity state.

\subsection{Structured sparse mixture-of-experts closure}
\label{subsec:sparse_moe_features}

The learned closure is conditioned on the local state, the projected
Galerkin contribution, the physical parameter, and a short temporal history.
Define
\begin{equation}
  \boldsymbol{g}^{(r)}_n
  =
  \mathscr{F}^{\mathrm{G}}_r
  \bigl(
    \boldsymbol{a}^{(r)}_n,
    \boldsymbol{b}^{(r)}_n;
    \mu
  \bigr),
  \label{eq:current_galerkin_feature}
\end{equation}
and let tildes denote regime-local standardization.  The current-state block
is
\begin{equation}
  \boldsymbol{\chi}^{(r)}_n
  =
  \operatorname{col}
  \left(
    \widetilde{\boldsymbol{a}}^{(r)}_n,
    \widetilde{\boldsymbol{b}}^{(r)}_n,
    \widetilde{\boldsymbol{g}}^{(r)}_n,
    \widetilde{\mu}
  \right),
  \label{eq:current_feature_block}
\end{equation}
and the two preceding states are encoded analogously,
\begin{equation}
  \boldsymbol{h}^{(r)}_{n,\ell}
  =
  \operatorname{col}
  \left(
    \widetilde{\boldsymbol{a}}^{(r)}_{n-\ell},
    \widetilde{\boldsymbol{b}}^{(r)}_{n-\ell},
    \widetilde{\boldsymbol{g}}^{(r)}_{n-\ell},
    \widetilde{\mu}
  \right),
  \qquad \ell=1,2.
  \label{eq:history_feature_block}
\end{equation}
The complete specialist feature vector is
\begin{equation}
  \boldsymbol{\xi}^{(r)}_n
  =
  \operatorname{col}
  \left(
    \boldsymbol{\chi}^{(r)}_n,
    \boldsymbol{h}^{(r)}_{n,1},
    \boldsymbol{h}^{(r)}_{n,2}
  \right).
  \label{eq:complete_specialist_features}
\end{equation}
The inclusion of $\boldsymbol{g}^{(r)}_n$ exposes the local projected vector
field directly to the learned closure, while the short history provides
information on recent temporal evolution that is not contained in the
instantaneous reduced state alone.

The closure is represented by a hierarchical sparse mixture of experts.  A
group router first selects a local expert group,
\begin{equation}
  g^{\star}
  =
  \underset{g\in\mathcal{G}}{\operatorname{arg\,max}}
  \;\pi_g^{\mathrm{grp}}(\boldsymbol{\xi}),
  \label{eq:group_top1}
\end{equation}
and, within the selected group, a channel-specific router activates a small
subset $\mathcal{T}^{c}(\boldsymbol{\xi};g^{\star})$ together with a shared
expert.  For channel $c\in\{u,p\}$,
\begin{equation}
  \boldsymbol{\mathcal{M}}^{c}_{\theta_r}(\boldsymbol{\xi})
  =
  \omega^{c}_{0}\,
  \mathcal{E}^{c}_{0,g^{\star}}(\boldsymbol{\xi})
  +
  \sum_{e\in\mathcal{T}^{c}(\boldsymbol{\xi};g^{\star})}
  \omega^{c}_{e}\,
  \mathcal{E}^{c}_{e,g^{\star}}(\boldsymbol{\xi}),
  \label{eq:channel_sparse_moe}
\end{equation}
with non-negative normalized weights,
\begin{equation}
  \omega^{c}_{e}\ge 0,
  \qquad
  \omega^{c}_{0}
  +\sum_{e\in\mathcal{T}^{c}}\omega^{c}_{e}=1.
  \label{eq:channel_sparse_weights}
\end{equation}
Separate channel routing allows the velocity closure and learned pressure map
to specialize differently while retaining a common reduced-state context.

Each expert has a structured decomposition
\begin{equation}
  \mathcal{E}^{c}_{e,g}(\boldsymbol{\xi})
  =
  \mathcal{N}^{c}_{e,g}(\boldsymbol{\xi})
  +W^{c,\mathrm{lin}}_{e,g}\boldsymbol{\xi}
  +\kappa W^{c,o}_{e,g}
  \left[
    \bigl({U^{c}_{e,g}}^{\mathsf T}\boldsymbol{\xi}\bigr)
    \odot
    \bigl({V^{c}_{e,g}}^{\mathsf T}\boldsymbol{\xi}\bigr)
  \right].
  \label{eq:physics_aware_expert}
\end{equation}
Here $\mathcal{N}^{c}_{e,g}$ is a nonlinear branch, the second term is an
affine contribution, and the third term is a low-rank bilinear interaction.
The matrices $U^{c}_{e,g},V^{c}_{e,g}\in\mathbb{R}^{m_r\times q_b}$ and
$W^{c,o}_{e,g}\in\mathbb{R}^{d_c\times q_b}$ define a rank-$q_b$ bilinear
map.  This branch provides an economical representation of pairwise
interactions among reduced and physical features and is consistent with the
low-order polynomial structure inherited from the projected convective
operator, without constraining the learned correction to reproduce the
Galerkin tensor.  Architectural hyperparameters such as $q_b$, the number of
routed experts, and network widths are specified with the numerical setup.

\subsection{Autonomous time advancement and rollout training}
\label{subsec:regime_specific_contracts}

Each specialist uses the same numerical form but an independent parameter set
$\theta_r$.  During a velocity macro-step, the pressure coefficients, recent
history, and physical parameter define a context $\mathcal{C}^{(r)}_n$.  The
stagewise right-hand side is
\begin{equation}
  F_r
  \bigl(\boldsymbol{a};\mathcal{C}^{(r)}_n\bigr)
  =
  \mathscr{F}^{\mathrm{G}}_r
  \bigl(\boldsymbol{a},\boldsymbol{b}^{(r)}_n;\mu\bigr)
  +
  \boldsymbol{\rho}^{u}_{\theta_r}
  \bigl(
    \boldsymbol{\xi}^{(r)}(\boldsymbol{a};\mathcal{C}^{(r)}_n)
  \bigr).
  \label{eq:frozen_context_rhs}
\end{equation}
With the classical fourth-order Runge--Kutta method,
\begin{align}
  \boldsymbol{k}_1
  &=F_r\bigl(\boldsymbol{a}^{(r)}_n;\mathcal{C}^{(r)}_n\bigr),\nonumber\\
  \boldsymbol{k}_2
  &=F_r\left(
      \boldsymbol{a}^{(r)}_n
      +\tfrac{\Delta t_r(\mu)}{2}\boldsymbol{k}_1;
      \mathcal{C}^{(r)}_n
    \right),\nonumber\\
  \boldsymbol{k}_3
  &=F_r\left(
      \boldsymbol{a}^{(r)}_n
      +\tfrac{\Delta t_r(\mu)}{2}\boldsymbol{k}_2;
      \mathcal{C}^{(r)}_n
    \right),\nonumber\\
  \boldsymbol{k}_4
  &=F_r\left(
      \boldsymbol{a}^{(r)}_n
      +\Delta t_r(\mu)\boldsymbol{k}_3;
      \mathcal{C}^{(r)}_n
    \right),\nonumber\\
  \boldsymbol{a}^{(r)}_{n+1}
  &=
  \boldsymbol{a}^{(r)}_n
  +\frac{\Delta t_r(\mu)}{6}
   \bigl(
     \boldsymbol{k}_1+2\boldsymbol{k}_2
     +2\boldsymbol{k}_3+\boldsymbol{k}_4
   \bigr).
  \label{eq:unified_rk4}
\end{align}
The updated velocity is then supplied to the algebraic pressure reconstruction
\eqref{eq:adaptive_pressure_closure}, after which the temporal history is
shifted and the autonomous rollout proceeds.

The specialist parameters are learned from closed-loop multi-step rollouts.
For a rollout of length $K$, the normalized reduced-coordinate objective is
\begin{equation}
  \mathcal{L}_{r}
  =
  \frac{1}{K}\sum_{k=1}^{K}
  \left[
    \frac{1}{d_u}
    \left\|
      \bigl(
        \boldsymbol{a}^{(r)}_{n+k}
        -\widehat{\boldsymbol{a}}^{(r)}_{n+k}
      \bigr)
      \oslash\boldsymbol{\sigma}^{(r)}_a
    \right\|_2^2
    +
    \frac{1}{d_p}
    \left\|
      \bigl(
        \boldsymbol{b}^{(r)}_{n+k}
        -\widehat{\boldsymbol{b}}^{(r)}_{n+k}
      \bigr)
      \oslash\boldsymbol{\sigma}^{(r)}_b
    \right\|_2^2
  \right].
  \label{eq:unified_rollout_loss}
\end{equation}
Training proceeds with progressively longer closed-loop horizons so that the
learned closure is optimized for recursive use rather than only one-step
approximation.  The schedule and optimizer settings are reported with the
numerical experiments.  The three specialists are trained independently;
inter-regime routing and fusion introduced in Section~\ref{sec:hierarchical_ral_moe}
do not modify their local reduced representations or dynamical parameters.
